% !TeX root = ../Challenges.tex
% !TeX spellcheck = en_US

\section{Open Research Challenge}
\label{ssec:orc}

The open research challenge is held to showcase innovations and research projects that support the development of the league beyond the core competition objectives. 

\subsection{Submission Requirements}
Teams must submit a brief description of their planned contribution by the day before the presentation at RoboCup 2026 to the Technical Committee (\url{hsl_tc@lists.robocup.org}).
Multiple contributions per team are allowed.
The description should include:
    \begin{itemize}
        \item \textbf{Objective}: The purpose and scope of the innovation.
        \item \textbf{Expected Impact}: How it benefits HSL or improves existing systems.
        \item \textbf{Implementation Feasibility}: Resources and time needed to accomplish/execute the contribution.
    \end{itemize}
    
\subsection{Presentation Format}
Participating teams have to prepare a short 10 minute oral presentation. The presentation must include a clear explanation of the contribution and its impacts. A live or video demonstration should be included if feasible.

Presented topics should focus on innovations beneficial to the league, but not necessarily related to winning matches. Examples include:

\begin{itemize}
\item Projects improving the safety and handling of robots
\item Projects that increase interoperability and collaboration
\item Projects that improve refereeing
\item Projects that enable game statistics and evaluation
\end{itemize}

Presentations are held publicly, encouraging members of all teams to join, even if they did not submit a contribution themselves.

\subsection{Evaluation Process}
The winner will be decided by a vote among the team leaders during the competition using the Condorcet method~\footnote{\url{https://en.wikipedia.org/wiki/Condorcet_method}}. Each participating team will vote for their top teams in order (excluding themselves). Teams are encouraged to evaluate the presentations based on the following criteria:

\begin{itemize}
\item \textbf{Relevance}: Alignment with HSL’s goals of growth and improvement.
\item \textbf{Novelty}: Originality and creativity of the idea.
\item \textbf{Feasibility}: Ease of application, practicality of implementation within league resources.
\item \textbf{Impact}: Potential benefit to the league or teams.
\item \textbf{Usability}: Ease of use, integration, and adaptability of the contribution.
\item \textbf{Live Demonstration}: Preference should be given to contributions that include a live demonstration.
\item \textbf{Documentation and Code Release}: Consideration of whether the code was made available beforehand for review and that it is well documented.
\end{itemize}

At a time decided by the designated referee, within one hour of the last demonstration if not otherwise specified, the captain of each team will submit the team’s rankings by filling out a form. Any points awarded by a team to itself will be disregarded.

\subsection{Full Code Release Requirement}
A \textbf{full code release} is mandatory for participating in this challenge, except for contributions that do not include any source code. The code must be publicly available by 2026-10-15
%no later than the \textbf{main code release deadline} from the HSL rule book 
and must be announced to the Technical Committee (\url{hsl_tc@lists.robocup.org}). Code must be well documented.
Ideally, the code should be released before the competition to allow interested teams and individuals to review and provide feedback!

\subsection{Publication}
The presentation slides, videos and all related material to the presented project must be made available to the Technical Committee on the day of the presentation. 
All presentations and their related code will be made publicly available on the HSL league website.